\documentclass[11pt,a4paper]{article}
\usepackage[margin=0.85in,top=1in,bottom=1in]{geometry}
\usepackage{amsmath,amssymb,amsfonts}
\usepackage{graphicx}
\usepackage{float}
\usepackage{enumitem}
\usepackage{microtype}
\usepackage{caption}
\usepackage[hidelinks]{hyperref}
\usepackage[most,breakable]{tcolorbox}
\tcbuselibrary{breakable,skins}
\usepackage{fontspec}
\setmainfont{DejaVu Serif}
\setsansfont{DejaVu Sans}
\setmonofont{DejaVu Sans Mono}
\captionsetup{font=small,labelfont=bf,skip=4pt}
\setlist{leftmargin=1.5em,topsep=2pt}

%% Breakable problem card — prevents content cutoff across pages
\newtcolorbox{solcard}[3]{
  breakable, enhanced,
  colback=white, colframe=gray!50, arc=2pt,
  title={\textbf{#1}\hfill\textsf{\small #3\,pts}},
  fonttitle=\normalsize\sffamily\bfseries,
  before upper={\small\textit{#2}\par\smallskip\hrule height 0.4pt\smallskip},
  top=4pt, bottom=6pt, left=7pt, right=7pt,
  parbox=false,
}

\title{\textbf{Binary thermodynamic cycles --- Solution}}
\author{\normalsize X23 (2023) — English translation}
\date{}
\begin{document}\maketitle\vspace{-0.5em}
\begin{solcard}{A1}{Express the amount of heat $Q_-$ removed from the gas per cycle through $T_1$, $T_2$, $T_4$, $R$, the amount of gas substance $\nu$ and the adiabatic exponent $\gamma$.}{0.70}

Solution in $TS$ coordinates.


The figure shows a qualitative diagram of the cycle under consideration in $TS$ coordinates.


Heat is removed from the gas in an isothermal process, so for $Q_-$ we get: $$Q_-=T_2(S_2-S_3) $$ where $S_2$ and $S_3$ - entropy gasа в точках $2$ and $3$ соответственно. Since изменение энтропии в адиабатических процессах equals нулю: $$S_2-S_3=S_1-S_4=\int\limits_{T_4}^{T_1}\cfrac{\nu C_pdT}{T}=\cfrac{\gamma\nu R}{\gamma-1}\ln\cfrac{T_1}{T_4} $$ from where we find


Solution in $pV$ or $pT$ coordinates.


The diagram shows a qualitative graph of the process in coordinates $pV$.


We use the Mendeleev-Clapeyron equation: $$pV=\nu RT $$ Тепло отводится от gasа в изотермическом процессе, в котором $\delta{Q}=\delta{A}$, therefore: $$\delta{Q}=pdV=\cfrac{\nu RTdV}{V}\Rightarrow Q_-=\nu RT_2\int\limits_{V_3}^{V_2}\cfrac{dV}{V}=\nu RT_2\ln\cfrac{V_2}{V_3} $$ Since процесс $23$ изотермический – $V_2/V_3=p_3/p _2$. Перепишем expression for $Q_-$: $$Q_-=\nu RT\ln\cfrac{p_3}{p_2} $$ For идеального gasа в квазистатическом адиабатическом процессе true equation Пуассона: $$p\sim{T^{\gamma/(\gamma-1)}} $$ therefore: $$\begin{cases} \cfrac{p_1}{p_2}=\left(\cfrac{T_1}{T_2}\right)^{\gamma/(\gamma-1)}\\ \cfrac{p_4}{p_3}=\left(\cfrac{T_4}{T_3}\right)^{\gamma/(\gamma-1)} \end{cases} \implies{\cfrac{p_3}{p_2}=\left(\cfrac{T_1}{T_4}\right)^{\gamma/(\gamma-1)}} $$ Since $T_3=T_2$ and $p_4=p_1$. Finally:

\par\smallskip\noindent\textbf{Answer:} \textbackslash{}textbf{Answer:} $$Q_-=\cfrac{\gamma \nu RT_2}{\gamma-1}\ln\cfrac{T_1}{T_4} $$
\begin{figure}[H]
  \centering
  \includegraphics[width=0.72\textwidth]{T2_S_EN_assets/fig1.jpg}
\end{figure}
\begin{figure}[H]
  \centering
  \includegraphics[width=0.72\textwidth]{T2_S_EN_assets/fig2.jpg}
\end{figure}
\end{solcard}

\begin{solcard}{A2}{Express the relation $T_1/T_4$ in terms of $\alpha$, $\beta$ and $\gamma$.}{0.80}

Let us write the adiabatic equation for the process $1-2$ in coordinates $pT$: $$p^{1-\gamma}T^\gamma=const\Rightarrow{\cfrac{T_1}{T_2}=\left(\cfrac{p_1}{p_2}\right)^{\frac{\gamma-1}{\gamma}}} $$ Since $p_1=p_4$ $$\cfrac{T_1}{T_2}=\beta^{\frac{\gamma-1}{\gamma}} $$ Also, since $T_2=T_4/\alpha$:

\par\smallskip\noindent\textbf{Answer:} \textbackslash{}textbf{Answer:} $$\cfrac{T_1}{T_4}=\cfrac{\beta^{\frac{\gamma-1}{\gamma}}}{\alpha} $$
\end{solcard}

\begin{solcard}{A3}{Obtain an expression for the cycle efficiency $\eta_1$. Express your answer in terms of $\alpha$, $\beta$ and $\gamma$. Find the numerical value of $\eta_1$ for $p_2=15~\text{kPa}$, $p_4=60~\text{kPa}$, $T_2=340~\text{K}$ and $T_4=400~\text{K}$.}{0.50}

Since $Q_+=C_p(T_1-T_4)$, we have: $$\eta_1=1-\cfrac{C_pT_2\ln(T_1/T_4)}{C_p(T_1-T_4)} $$ whence


For $\alpha$ and $\beta$ we have: $$\alpha=\cfrac{20}{17}\qquad \beta=4 $$ from where:

\par\smallskip\noindent\textbf{Answer:} \textbackslash{}textbf{Answer:} $$\eta_1=1-\cfrac{\ln\left(\beta^{\frac{\gamma-1}{\gamma}}/\alpha\right)}{\beta^{\frac{\gamma-1}{\gamma}}-\alpha} $$
\end{solcard}

\begin{solcard}{B1}{Find the ratio of the masses of mercury and water $m_\text{(р)}/m_\text{(в)}$ participating in the cycle.}{0.50}

Solution in $TS$ coordinates.


The figure shows qualitative diagrams of the considered cycles in $TS$ coordinates.


Since the amounts of heat received and given up by water and mercury vapor, respectively, in isobaric polytropic processes are the same at the same temperature changes, the heat capacities in isobaric processes are also the same: $$C_{p\text{(в)}}=\cfrac{m_\text{(в)}\gamma_\text{(в)}R}{\mu_\text{(в)}(\gamma_\text{(в)}-1)}=\cfrac{m_\text{(р)}\gamma_\text{(р)}R}{\mu_\text{(р)}(\gamma_\text{(р)}-1)}=C_{p\text{(р)}} $$ from where:

\par\smallskip\noindent\textbf{Answer:} \textbackslash{}textbf{Answer:} $$\cfrac{m_\text{(р)}}{m_\text{(в)}}=\cfrac{\mu_\text{(р)}}{\mu_\text{(в)}}\cfrac{\gamma_\text{(р)}-1}{\gamma_\text{(р)}}\cfrac{\gamma_\text{(в)}}{\gamma_\text{(в)}-1}=17{,}81 $$
\begin{figure}[H]
  \centering
  \includegraphics[width=0.72\textwidth]{T2_S_EN_assets/fig3.jpg}
\end{figure}
\end{solcard}

\begin{solcard}{B2}{Using graphs of the pressure of saturated vapors of mercury and water versus temperature, estimate the pressure values ​​$p_{c\text{(р)}}$ and $p_{c\text{(в)}}$.}{0.40}

The required pressures are equal to the pressures of saturated vapors at temperatures $t_4$ for mercury and $t_2$ for water. From the graphs we find:

\par\smallskip\noindent\textbf{Answer:} \textbackslash{}textbf{Answer:} $$p_{c\text{(р)}}=6~\text{kPa} $$
\end{solcard}

\begin{solcard}{B3}{Find the mercury and water cycle efficiencies $\eta_\text{(р)}$ and $\eta_\text{(в)}$ respectively.}{0.80}

Note: in all paragraphs below it is assumed that $T_i=t_i+T_0$, where $T_0=273~\text{K}$. The change in entropy in isothermal and isobaric processes is the same, therefore: $$\Delta{S}=\int\limits_{T_4}^{T_3}\cfrac{C_pdT}{T}=C_p\ln\cfrac{T_3}{T_4} $$ Heat, поrayенное ртутью в изотермическом процессе, equals: $$Q_{+\text{(р)}}=T_1\Delta{S} $$ Heat, отданное водяным vaporом в изотермическом процессе, equals: $$Q_{-\text{(в)}}=T_2\Delta{S} $$ В изобарном процессе водяной vapor поrayает а ртуть отgives heat $Q_p$, equalsе: $$Q_p=C_p(T_3-T_4) $$ therefore answerы for $\eta_\text{(р)}$ and $\eta_\text{(в)}$ are as follows:


The change in entropy in isothermal and isobaric processes is the same, therefore: $$\Delta{S}=\int\limits_{T_4}^{T_3}\cfrac{C_pdT}{T}=C_p\ln\cfrac{T_3}{T_4} $$ Heat, поrayенное ртутью в изотермическом процессе, equals: $$Q_{+\text{(р)}}=T_1\Delta{S} $$ Heat, отданное водяным vaporом в изотермическом процессе, equals: $$Q_{-\text{(в)}}=T_2\Delta{S} $$ В изобарном процессе водяной vapor поrayает а ртуть отgives heat $Q_p$, equalsе: $$Q_p=C_p(T_3-T_4) $$ therefore answerы for $\eta_\text{(р)}$ and $\eta_\text{(в)}$ the following:

\par\smallskip\noindent\textbf{Answer:} \textbackslash{}textbf{Answer:} $$\eta_\text{(р)}=1-\cfrac{T_3-T_4}{T_1\ln(T_3/T_4)}=0{,}326 $$
\end{solcard}

\begin{solcard}{B4}{Find the efficiency $\eta_\text{gп}$ of the binary gas-steam-water cycle. Compare it with $\eta_\text{(р)}$ and $\eta_\text{(в)}$.}{0.80}

Solution in $TS$ coordinates.


Since in an isobaric process gases exchange heat with each other, heat is supplied from the outside to the system participating in the binary cycle only in the mercury isothermal process. Then the efficiency of the binary cycle is the same as that of the Carnot cycle, consisting of two isotherms and two adiabats with temperatures $T_1$ and $T_2$: $$\eta_\text{Kарно}=1-\cfrac{T_\text{х}}{T_\text{н}} $$ Finally:

\par\smallskip\noindent\textbf{Answer:} \textbackslash{}textbf{Answer:} $$\eta_\text{gп}=1-\cfrac{T_2}{T_1}=0{,}657>\eta_\text{(р)}{,}\eta_\text{(в)} $$
\end{solcard}

\begin{solcard}{C1}{Find the temperatures of the system at points $2$, $3$, $4$, $5$ and $6$.}{0.40}

Since the process $23$ is isothermal, and the change in temperature in section $34$ can be neglected by condition, the temperatures $t_2$, $t_3$ and $t_4$ are equal to the temperature of saturated steam at pressure $p_3=16~\text{kPa}$, i.e.:


The system temperatures at points $5$ and $6$ correspond to the temperature of saturated steam at pressure $p_4=2{,}8~\text{MPa}$, i.e.:

\par\smallskip\noindent\textbf{Answer:} \textbackslash{}textbf{Answer:} $$t_2=t_3=t_4=t'_1=55^{\circ}\mathrm{C} $$
\end{solcard}

\begin{solcard}{C2}{Find the amount of heat $Q_+$ supplied to the system per cycle.}{0.60}

The figure shows a qualitative diagram of the cycle under consideration in $TS$ coordinates.


In section $4-5$ the liquid heats up, in section $5-6$ it evaporates, and section $6-1$ corresponds to the overheating of water vapor. Then for $Q_+$ we have: $$Q_+=cm(t_5-t_4)+L(t_5)m+\cfrac{\gamma_\text{(в)}mR(t_1-t_6)}{\mu_{\text{(в)}}(\gamma_\text{(в)}-1)} $$ Since $t_6>180^{\circ}\mathrm{C}$ - adiabatic index водяного vaporа equals $9/7$. Substituting numerical values, we find:

\par\smallskip\noindent\textbf{Answer:} \textbackslash{}textbf{Answer:} $$Q_+=57{,}35~\text{MJ} $$
\begin{figure}[H]
  \centering
  \includegraphics[width=0.72\textwidth]{T2_S_EN_assets/fig4.jpg}
\end{figure}
\end{solcard}

\begin{solcard}{C3}{Find the amount of heat $Q_-$ removed from the system per cycle.}{0.80}

The amount of heat given off per cycle in an isothermal process is equal to: $$Q_-=T_2(S_2-S_3) $$ Let us find изменение энтропии другим способом, taking into account that процессы $1-2$ and $3-4$ являются адиабатными: $$S_2-S_3=S_5-S_4+S_6-S_5+S_1-S_6=m\left(c\ln\cfrac{T_5}{T_4}+\cfrac{L(t_5)}{T_5}+\cfrac{\gamma_\text{(в)}R}{\mu_\text{(в)}(\gamma_\text{(в)}-1)}\ln\cfrac{T_1}{T_6}\right) $$ Thus: $$Q_-=mT_2\left(c\ln\cfrac{T_5}{T_4}+\cfrac{L(t_5)}{T_5}+\cfrac{\gamma_\text{(в)}R}{\mu_\text{(в)}(\gamma_\text{(в)}-1)}\ln\cfrac{T_1}{T_6}\right) $$ Substituting the numerical values, we find:

\par\smallskip\noindent\textbf{Answer:} \textbackslash{}textbf{Answer:} $$Q_-=36{,}83~\text{MJ} $$
\end{solcard}

\begin{solcard}{C4}{Find the gas work $A$ per cycle and its efficiency $\eta_2$.}{0.20}

Let's use the law of conservation of energy: $$A=Q_+-Q_- $$ from where:


By definition: $$\eta_2=\cfrac{A}{Q_+}=1-\cfrac{Q_-}{Q_+} $$ from:

\par\smallskip\noindent\textbf{Answer:} \textbackslash{}textbf{Answer:} $$A=20{,}52~\text{MJ} $$
\end{solcard}

\begin{solcard}{D1}{Find the temperatures at points $b_\text{в}$, $c_\text{в}$, $d_\text{в}$, $e_\text{в}$, $f_\text{в}$ of the water cycle.}{0.20}

The temperatures at points $b_{(\text{в})}$, $c_{(\text{в})}$ and $d_{(\text{в})}$ are the same and equal to the temperature of saturated water vapor at pressure $p_{c(\text{в})}$:


The temperatures at points $e_{(\text{в})}$ and $f_{(\text{в})}$ are the same and equal to the temperature of saturated water vapor at pressure $p_{e(\text{в})}$:

\par\smallskip\noindent\textbf{Answer:} \textbackslash{}textbf{Answer:} $$t_{b(\text{в})}=t_{c(\text{в})}=t_{d(\text{в})}=t_1=96^{\circ}\mathrm{C} $$
\end{solcard}

\begin{solcard}{D2}{Find the temperature at point $a_\text{в}$ of the water cycle.}{1.00}

The figure shows qualitative diagrams of the considered cycles in $TS$ coordinates.


Since the processes $a_\text{в}-b_\text{в}$ and $c_\text{в}-d_\text{в}$ are adiabatic, the change in entropy in sections $d_\text{в}-e_\text{в}-f_\text{в}-a_\text{в}$ is equal to the change in entropy in section $c_\text{в}-b_\text{в}$: $$S_{b\text{(в)}}-S_{c\text{(в)}}=S_{e\text{(в)}}-S_{d\text{(в)}}+S_{f\text{(в)}}-S_{e\text{(в)}}+S_{a\text{(в)}}-S_{f\text{(в)}} $$ For изменении энтропии воды на участках $d_\text{в}-e_\text{в}-f_\text{в}-a_\text{в}$ we have: $$S_{a\text{(в)}}-S_{d\text{(в)}}=m_\text{в}\left(c_\text{в}\ln\cfrac{T_2}{T_1}+\cfrac{L_\text{в}(T_2)}{T_2}+\cfrac{\gamma_\text{в}R}{\mu_\text{в}(\gamma_\text{в}-1)}\ln\cfrac{T_{a\text{(в)}}}{T_2}\right) $$ For изменения энтропии на участке $c_\text{в}-b_\text{в}$ we have: $$S_{b\text{(в)}}-S_{c\text{(в)}}=\cfrac{m_\text{в}L_\text{в}(T_1)}{T_1}+\cfrac{A_\text{пар}}{T_1} $⟦11 ⟧A_\text{par}$ - работа водяного пара при расширении. Найдём её: $$A_\textbackslash{}text{vapor}=\textbackslash{}int\textbackslash{}limits_{V_0}^{V_\textbackslash{}text{к}}pdV=\textbackslash{}nu RT\textbackslash{}int\textbackslash{}limits_{V_0}^{V_\textbackslash{}text{к}}\textbackslash{}cfrac{dV}{V}=\textbackslash{}nu RT\textbackslash{}ln\textbackslash{}cfrac{V_\textbackslash{}text{к}}{V_0} $$ Поскольку пар расширяется по изотерме: $$\textbackslash{}cfrac{V_\textbackslash{}text{к}}{V_0}=\textbackslash{}cfrac{p_0}{p_\textbackslash{}text{к}}=\textbackslash{}cfrac{p_{c\textbackslash{}text{(в)}}}{p_{b\textbackslash{}text{(в)}}} $$ Мы получили уравнение, содержащее $T_a$: $$\textbackslash{}cfrac{L_\textbackslash{}text{в}(T_1)}{T_1}+\textbackslash{}cfrac{R}{\textbackslash{}mu_\textbackslash{}text{в}}\textbackslash{}ln\textbackslash{}cfrac{p_{c\textbackslash{}text{(в)}}}{p_{b\textbackslash{}text{(в)}}}=c_\textbackslash{}text{в}\textbackslash{}ln\textbackslash{}cfrac{T_2}{T_1}+\textbackslash{}cfrac{L_\textbackslash{}text{в}(T_2)}{T_2}+\textbackslash{}cfrac{\textbackslash{}gamma_\textbackslash{}text{в}R}{\textbackslash{}mu_\textbackslash{}text{в}(\textbackslash{}gamma_\textbackslash{}text{в}-1)}\textbackslash{}ln\textbackslash{}cfrac{T_{a\textbackslash{}text{(в)}}}{T_2} $⟦ 19⟧L_\text{в}(T_1)$ and $L_\text{в}(T_2)$: $$L_\text{в}(T_1)=2{,}27~\text{MJ}\qquad L_\text{в}(T_2)=1{,}98~\text{MJ} $$ Finally we get:

\par\smallskip\noindent\textbf{Answer:} \textbackslash{}textbf{Answer:} $$t_{a\text{(в)}}\approx{677^{\circ}\mathrm{C}} $$
\begin{figure}[H]
  \centering
  \includegraphics[width=0.72\textwidth]{T2_S_EN_assets/fig5.jpg}
\end{figure}
\end{solcard}

\begin{solcard}{D3}{Find the mercury pressure at section $d_\text{р}-e_\text{р}-a_\text{р}$.}{1.30}

Let us denote the temperature at point $e_\text{р}$ as $T_4$. Then for the change in entropy in section $d_\text{р}-e_\text{р}-a_\text{р}$ we have: $$S_{a\text{(р)}}-S_{d\text{(р)}}=m_\text{р}\left(c_\text{р}\ln\cfrac{T_4}{T_2}+\cfrac{L_\text{р}(T_4)}{T_4}\right) $$ On the other hand, это же изменение энтропии equals изменению энтропии на участке $c_\text{р}-b_\text{р}$, которое, в свою очередь, equals изменению энтропии воды на участке $e_\text{в}-f_\text{в}$: $$S_{f\text{(в)}}-S_{f\text{(в)}}=\cfrac{m_\text{в}L_\text{в}(T_2)}{T_2} $$ For/Atравнивая, we get зависимость $L_\text {р}(T_4)$: $$L_\text{р}(T_4)=\cfrac{m_\text{в}L_\text{в}(T_2)T_4}{m_\text{р}T_2}-c_\text{р}T_4\ln\cfrac{T_2}{T_4} $$ Построим данную функцию на границе зависимости $L_\text{р}(t)$ and let us find solution: $$t_4\approx 610^{\circ}\mathrm{C} $$ Ссу we determine $L_\text{р}(t_4)$: $$L_\text{р}(t_4)\approx 288~\text{kJ}/\text{kg} $$ The pressure of mercury in the desired area is equal to the pressure of saturated vapor of mercury at a given temperature. Finally we have:

\par\smallskip\noindent\textbf{Answer:} \textbackslash{}textbf{Answer:} $$p_{e\text{(р)}}\approx 2{,}6~\text{MPa} $$
\end{solcard}

\begin{solcard}{D4}{Find the efficiency of the binary Rankine cycle with mercury and water $\eta_3$.}{1.00}

Mercury in section $b_\text{р}-c_\text{р}$ exchanges heat with water in section $e_\text{в}-f_\text{в}$, so heat is supplied from the outside to mercury in section $d_\text{р}-e_\text{р}-a_\text{р}$, as well as to water in sections $d_\text{в}-e_\text{в}$ and $f_\text{в}-a_\text{в}$. Then for the amount of heat supplied per unit mass of water $Q_{+3}/m_\text{в}$ we obtain: $$\cfrac{Q_{+3}}{m_\text{в}}=\left(c_\text{в}(t_2-t_1)+\cfrac{\gamma_\text{в}R(t_3-t_2)}{\mu_\text{в}(\gamma_\text{в}-1)}\right)+\cfrac{m_\text{р}}{m_\text{в}}\left(c_\text{р}(t_4-t_2)+L_\text{р}(t_4)\right)\approx{4{,}976~\cfrac{\text{MJ}}{\text{kg}}} $$ Тепло в системе отводится только от воды на участке $b_\text{в}-c_\text{в}$, therefore for отведённого количества теплоты на единицу массы воды $Q_{-3}/m_\text{в}$ we have: $$\cfrac{Q_{-3}}{m_\text{в}}=\cfrac{T_1(S_{b\text{(в)}}-S_{c\text{(в)}})}{m_\text{в}} $$ From resultов partа $\mathrm{D2}$: $$\cfrac{S_{b\text{в}}-S_{c\text{в}}}{m_\text{в}}=\left(c_\text{в}\ln\cfrac{T_2}{T_1}+\cfrac{L_\text{в}(T_2)}{T_2}+\cfrac{\gamma_\text{в}R}{\mu_\text{в}(\gamma_\text{в}-1)}\ln\cfrac{T_{a\text{(в)}}}{T_2}\right) $$ Hence: $$\cfrac{Q_{-3}}{m_\text{в}}=T_1\left(c_\text{в}\ln\cfrac{T_2}{T_1}+\cfrac{L_\text{в}(T_2)}{T_2}+\cfrac{\gamma_\text{в}R}{\mu_\text{в}(\gamma_\text{в}-1)}\ln\cfrac{T_{a\text{(в)}}}{T_2}\right)\approx{2{,}484~\cfrac{\text{MJ}}{Кг}} $$ Final:

\par\smallskip\noindent\textbf{Answer:} \textbackslash{}textbf{Answer:} $$\eta_3\approx{0{,}50} $$
\end{solcard}

\end{document}